% TOP_PROBLEM: {"problemId": "problem.consistency-of-reinhardt-cardinals-without-choice", "canonicalTitle": "Consistency of Reinhardt Cardinals without Choice", "releaseRank": 152, "primaryDomain": "logic_foundations_set_theory", "primaryDomainLabel": "Logic, foundations & set theory", "displayStatus": "open", "releaseStatus": "open"}
% !TeX program = lualatex
% Definitions and sources retained from research_batch_151-153.tex; research is in GPT_6_Astra_Ultra.
\documentclass[11pt]{article}
\usepackage[margin=25mm]{geometry}
\usepackage{fontspec}
\setmainfont{Latin Modern Roman}
\usepackage{amsmath,amssymb,mathtools}
\usepackage{unicode-math}
\setmathfont{Latin Modern Math}
\usepackage{enumitem,needspace}
\usepackage{xurl,hyperref}
\hypersetup{colorlinks=true,urlcolor=blue}
\setcounter{secnumdepth}{1}
\setlength{\parindent}{0pt}
\setlength{\parskip}{5pt}
\setlength{\emergencystretch}{3em}
\setlist[itemize]{leftmargin=3em,itemsep=5pt,topsep=4pt,parsep=0pt}
\allowdisplaybreaks
\newcommand{\N}{\mathbb{N}}
\newcommand{\Z}{\mathbb{Z}}
\newcommand{\Q}{\mathbb{Q}}
\newcommand{\R}{\mathbb{R}}
\newcommand{\C}{\mathbb{C}}
\newcommand{\F}{\mathbb{F}}
\newcommand{\A}{\mathbb{A}}
\newcommand{\PP}{\mathbb P}
\newcommand{\E}{\mathbb{E}}
\newcommand{\eps}{\varepsilon}
\newcommand{\dd}{\,\mathrm d}
\newcommand{\sref}[2]{\hyperlink{source:#1:#2}{[S#2]}}
\newcommand{\eref}[2]{\hyperlink{extra:#1:#2}{[E#2]}}
% Turn the next switch off for a shorter reading copy. The quotations preserve
% every exactTarget field, including qualifications omitted from a short summary.
\newif\ifshowcatalogue
\showcataloguetrue
\newcommand{\cataloguescope}[1]{\ifshowcatalogue
 \paragraph{Exact scope in the supplied catalogue.}
 \begin{quote}\small #1\end{quote}\fi}
\begin{document}
\setcounter{section}{151}
% ================================================================
% problemId: problem.consistency-of-reinhardt-cardinals-without-choice
% source releaseRank: 152; source releaseStatus: open
% exactTarget SHA-256: 94e7bff9ce564568770d0988780c6b02f381ee100e2c52c8bf74079f3b859e3d
\clearpage
\section{\texorpdfstring{Consistency of Reinhardt Cardinals without Choice}{Consistency of Reinhardt Cardinals without Choice}}\label{152}
\subsection{Definitions and mathematical statement}
Use the first-order language of set theory enlarged by a unary function symbol $j$. The theory $\mathrm{ZF}(j)$ includes the indicated ZF axiom schemes in that expanded language; this choice matters because formulas may mention $j$. Require $j$ to be nontrivial and $\Sigma_1$-elementary:
\[
 \varphi(a_1,\ldots,a_m)\ \longleftrightarrow\
 \varphi(j(a_1),\ldots,j(a_m))
\]
for each $\Sigma_1$ formula in the membership language. The theory specified by the source derives the full elementarity scheme. Its critical point, the least ordinal moved by $j$, is a Reinhardt cardinal. The target is the consistency of this exact theory, denoted ZFR, not a theory with the axiom of choice added. Consistency means that no contradiction is derivable by a finite formal proof; a relative-consistency result must state its additional assumptions.
\par\smallskip\textit{Formulation and concept references:} \sref{152}{1}.
\cataloguescope{Determine whether ZFR is consistent, where ZFR is the first-order theory ZF(j) plus the assertion that a nontrivial function symbol j:V->V is Sigma\_1-elementary and hence fully elementary as a theorem scheme.}
\subsection{Short English statement}
Can set theory without the axiom of choice consistently admit a nontrivial elementary self-embedding of the entire universe, with the axiom schemes specified in ZFR?
\subsection{Sources}
\begin{itemize}\raggedright
\item[\textnormal{[S1]}]\hypertarget{source:152:1}{} F. Schlutzenberg, Reinhardt cardinals and iterates of V, Ann. Pure Appl. Logic 173 (2022).
\url{https://arxiv.org/abs/2002.01215}.
{\small\textit{Catalogue formulation source.}}
\item[\textnormal{[E1]}]\hypertarget{extra:152:1}{} Marwan Salam Mohammd,
\emph{Reinhardt cardinals and eventually dominating functions},
Archive for Mathematical Logic 65 (2026), 691--699.
\url{https://doi.org/10.1007/s00153-026-01018-2}.
Author manuscript: \url{https://arxiv.org/html/2505.00637v1} (1 May 2025),
Theorem 2.2 and \S\S3,5.
{\small\textit{Manuscript consulted and publication metadata checked
24 September 2026; the Choice-dependent argument in \S5 is not applied to ZFR.}}
% END RESEARCH ADDITION REFERENCES-152
\end{itemize}

\end{document}
